\documentclass[book.tex]{subfiles}

\begin{document}
\chapter{Physics Informed Neural Networks}

\label{chap:pinns}
\noindent\rule{11cm}{0.4pt} \\
\textbf{Prerequisites:}  \autoref{chap:pde}, \autoref{chap:numerical_pde}, \autoref{chap:univ_approx}  \\
\textbf{Difficulty Level:} ** \\
\noindent\rule{11cm}{0.4pt}
\vspace{5mm}

Physics-informed neural networks (PINNs) are a family of neural networks trained to solve supervised learning tasks by using any given law of physics, typically described by general nonlinear partial differential equations (PDEs)\cite{raissi2019physics}, as a regularizer. PINNs leverage the fact that deep neural networks are universal function approximators to solve nonlinear problems without needing to commit to prior assumptions such as linearization, or local time-stepping. PINNs differentiate neural networks with respect to their input coordinates and model parameters, and can be constrained to respect any symmetries, invariances, or conservation principles originating from the physical laws that govern the observed data. This simple yet powerful construction allows us to construct a new class of numerical solvers for partial differential equations, as well as new data-driven approaches for model inversion and systems identification.



\section{PINNs Formulation}
Consider a parametrized and nonlinear partial differential equations of the following general form 
\begin{align}
    u_t + \mathcal{N}[u;\theta]=0,\quad x \in \Omega, t\in [0,T]
    \label{eq_generalform}
\end{align}
where $u(t, x)$ denotes the hidden solution, $\mathcal{N[\cdot;\theta]}$ is a nonlinear operator parametrized by $\theta$, and $\Omega$ is a subset of $\mathbb{R}^D$. This formulation applies to a wide range of problems in physical systems including conservation laws, diffusion processes, kinetic equations and physical systems.

We can speak of two methods of solving PINNs.
\begin{enumerate}
    \item Given fixed model parameters $\theta$, what can be said about the unknown hidden state $u(t, x)$ of the system?
    \item Can we determine the parameters $\theta$ that best describe observed data?
\end{enumerate}

The construction and training of PINNs can be done through two distinct algorithms based on the specific application area and the nature of the associated ground truth data:
\begin{enumerate}
    \item Continuous-time 
    \item Discrete-time 
\end{enumerate} 
We will learn about these two approaches in the following section. 
\section{Continuous Time Approximation}

The approach below defines an algorithm that can be used to construct a continuous time model. Let's define $f(t, x)$ to be given by the equation on the left hand side,

\begin{align}
    f &=u_t + \mathcal{N}[u; \theta]
    \label{eq_contTime}
\end{align}

Our approach is to approximate $u(t,x)$ by a deep neural network. This assumption along with equation \ref{eq_contTime} gives rise to the physics-informed neural network $f(t,x)$. The network can be derived by applying automatic differention to compute $\mathcal{N}[u; \theta]$. The shared parameters $\theta$ between the neural networks $u(t, x)$ and $f(t, x)$ can be learned by minimizing the mean-squared error loss

\begin{align}
    \mathcal{L}_{MSE}=\textrm{MSE}_u + \textrm{MSE}_f 
\end{align}

where 
\begin{align}
    \textrm{MSE}_u &=\frac{1}{N} \sum_{i=1}^{N}|u(t^i,x^i)-u^i|^2\\
    \textrm{MSE}_f &=\frac{1}{N} \sum_{i=1}^{N}|f(t^i,x^i)|^2
\end{align}
Above, the initial and boundary training data on $u(t, x)$ are denoted by 
\begin{align}
\{{t^i, x^i, u^i}\}^{N}_{i=1}
\end{align} 

The loss function $\textrm{MSE}_u$ penalizes deviations from the initial and boundary data while $\textrm{MSE}_f$ enforces the structure imposed by the general equation (equation \ref{eq_generalform}) at a finite set of collocation points. 

We don't need to optimize $u$ and $f$ at the same points so we can also generalize this loss function to 
\begin{align}
    \textrm{MSE}_u &=\frac{1}{N_u} \sum_{i=1}^{N_u}|u(t_u^i,x_u^i)-u^i|^2\\
    \textrm{MSE}_f &=\frac{1}{N_f} \sum_{i=1}^{N_f}|f(t_f^i,x_f^i)|^2
\end{align}

Here initial and boundary training data on $u(t, x)$ are denoted by 
\begin{align}
\{{t^i_u, x^i_u, u^i}\}^{N_u}_{i=1}, \quad \{{t^i_f, x^i_f}\}^{N_f}_{i=1}
\end{align} 
specify the collocation points for $f(t, x)$. 

\begin{lstlisting}[language=python]

class PINN(equation: Expression, u: Module):

  $\mathcal{N}$ = tensorize(equation)
  f = u + $\mathcal{N}$(u)

  def $\lambda$(x, t):
    return u(x, t)
    
  def loss(X, U):
    return sum($\|$u(X[i]) - U[i]$\|^2$, i) + sum($\|$f(X[i])$\|^2$, i)

\end{lstlisting}


\section{Discrete-Time PINNs}
The approach discussed above approximates the solution $u$ across a continuous range of time points. For some problems, it can be useful to discretize time instead. The general form of Runge-Kutta methods (\autoref{chap:runge_kutta}) with $q$ stages is applied to the governing equation to arrive at the following:

\begin{align}
u^{n+c_i}&=u^n-\Delta t\sum_{i=1}^{q}a_{ij}\mathcal{N}[u^{n+c_j}; \theta], \quad i=1,\cdot\cdot\cdot,q,\\
u^{n+1}&=u^n-\Delta t\sum_{i=1}^{q}b_{j}\mathcal{N}[u^{n+c_j}; \theta].
\end{align}

Here $a_{ij}$ and $b_j$ are the Runge-Kutta coefficients for $q$ stages. In the formulation, 
\begin{align}
u^{n+c_j}(x) &=u(t^n+c_j\Delta t,x)
\end{align}
for $j=1,\cdot\cdot\cdot,q$. The above equations can also be expressed as

\begin{align}
u^n &=u^n_i,\quad i=1,\cdot\cdot\cdot,q,\\
u^n &=u^n_{q+1}
\end{align}

where
\begin{align}
u_i^n&:=u^{n+c_i}+\Delta t\sum_{j=1}^{q}a_{ij}\mathcal{N}[u^{n+c_j}], \quad i=1,\cdot\cdot\cdot,q,\\
u^{n}_{q+1}&=u^{n+1}+\Delta t\sum_{i=1}^{q}b_{j}\mathcal{N}[u^{n+c_j}].
\end{align}

We can place a multi-output neural network prior on
\begin{align}
\label{eq_discrete}
    [u^{n+c_1}(x),\cdot\cdot\cdot,u^{n+c_q}(x),u^{n+1}(x)]
\end{align}
The assumption along with equations \ref{eq_discrete} result in a physics-informed neural network that takes $x$ as an input and gives as the output 
\begin{align}
    [u^{n}_1(x),\cdot\cdot\cdot,u^{n}_q(x),u_{q+1}^{n+1}(x)]
\end{align}

The discrete time network is trained by minimizing the following sum of squares loss (SSE)

\begin{align}
    \mathcal{L}_{SSE} &= \textrm{SSE}_n + \textrm{SSE}_b
\end{align}
where 
\begin{align}
    \textrm{SSE}_n &= \sum_{j=1}^{q+1} \sum_{i=1}^{N_n} |u^n_j(x^{n,i}) - u^{n,i}|^2 \\
    \textrm{SSE}_b &= \sum_{i=1}^q f(u^{n+c_i}) + f(u^{n+1}) \\
\end{align}
where $f$ is the general nonlinear constraint. In the above equations, the dataset at time $t^n$ is given by
\begin{align}
    \{x^{n,i}, u^{n,i}\}_{i=1}^{N_n}
\end{align}

\section{Exercises}
\begin{enumerate}
    \item Write down the explicit continuous-time PINN loss function for Schrodinger's equation in 1 dimension.
\end{enumerate}

\end{document}